%=====================================================================
% Auswertung endlicher Bitfolgen
% Kompilation: xelatex -> bibtex -> xelatex -> xelatex
% Literatur:   references.bib (biblatex, backend=bibtex)
%=====================================================================
\documentclass[11pt,a4paper]{scrartcl}

%-----------------------------------------------------
% Schrift (XeLaTeX). Fallback auf Standardschriften,
% falls die Nennschriften nicht installiert sind.
%-----------------------------------------------------
\usepackage{fontspec}
\IfFontExistsTF{EB Garamond}{\setmainfont[Ligatures=TeX]{EB Garamond}}{}
\IfFontExistsTF{Latin Modern Sans}{\setsansfont[Ligatures=TeX]{Latin Modern Sans}}{}
\IfFontExistsTF{Latin Modern Mono}{\setmonofont[Ligatures=TeX]{Latin Modern Mono}}{}

%-----------------------------------------------------
% Pakete
%-----------------------------------------------------
\usepackage[left=3cm,right=3.5cm,top=3cm,bottom=3.5cm]{geometry}
\usepackage[ngerman]{babel}
\usepackage{amsmath,amssymb,amsthm}
\usepackage{graphicx}
\usepackage{booktabs}
\usepackage[normalem]{ulem}   % fuer \xout in Tabelle 2
\usepackage{xcolor}
\usepackage{listings}
\usepackage[german=quotes]{csquotes}
\usepackage{microtype}
\usepackage[backend=bibtex,style=numeric,maxbibnames=9]{biblatex}
\addbibresource{references.bib}
\usepackage[hidelinks,pdfusetitle]{hyperref}

%-----------------------------------------------------
% Pseudocode-Listings
%-----------------------------------------------------
\lstset{%
  basicstyle=\ttfamily\small,
  frame=single,
  captionpos=b,
  breaklines=true,
  keepspaces=true,
  tabsize=2,
  showstringspaces=false
}
\renewcommand{\lstlistingname}{Algorithmus}

%-----------------------------------------------------
% Theorem-Apparat (durchgaengig gemeinsamer Zaehler)
%-----------------------------------------------------
\theoremstyle{plain}
\newtheorem{satz}{Satz}
\newtheorem{lemma}[satz]{Lemma}
\newtheorem{korollar}[satz]{Korollar}
\theoremstyle{definition}
\newtheorem{defi}[satz]{Definition}
\newtheorem{beispiel}[satz]{Beispiel}
\theoremstyle{remark}
\newtheorem{bem}[satz]{Bemerkung}

%-----------------------------------------------------
% Titel
%-----------------------------------------------------
\subject{Ein Testverfahren und die konstruktive Erzeugung von
  Bitfolgen mit exponentiell kumulierten Häufigkeitsverteilungen ihrer Läufe}
\title{Auswertung endlicher Bitfolgen}
\author{Nicolas Windt}
\date{} % TODO(Meta): bei Einreichung festes Datum setzen

\begin{document}

\maketitle

%-----------------------------------------------------
% Abstract
%-----------------------------------------------------
\begin{abstract}
Wir führen den \emph{Einschritttest} ein, ein Auswertungsverfahren für
endliche Bitfolgen, das die Häufigkeiten der Läufe und Teilläufe einer Folge
kumuliert. Die Anwendung auf die Ausgabe des Zufallsgeräts \texttt{/dev/random}
einer Linux-Installation zeigt einen annähernd exponentiell fallenden Verlauf
der kumulierten Häufigkeiten, im Einklang mit der geometrischen Verteilung der
Lauflängen unabhängiger Bits. Motiviert durch diese Beobachtung untersuchen
wir Bitfolgen, deren kumulierte Häufigkeitsverteilung exakt exponentiell
fällt, und fassen sie zu Klassen $B_N$ zusammen. Wir zeigen, dass die
Verteilung der Läufe ohne Teilläufe (Grundverteilung) der Elemente aus $B_N$
einer strengen Rekursion genügt, und leiten eine Identität für die Anzahl der
abgegoltenen Teilläufe her. Darauf aufbauend geben wir einen iterativen und
einen rekursiven Algorithmus zur Erzeugung eines Elements aus $B_N$ an.
Schließlich ordnen wir jeder natürlichen Zahl eine eindeutige Zusammensetzung
paarweise verschiedener Klassen $B_N$ zu, erhalten daraus eine Familie
zusammengesetzter Häufigkeitsverteilungen und diskutieren den dadurch
induzierten Zufallsbegriff.
\end{abstract}

\vspace{0.5em}
\noindent\textbf{Schlüsselwörter:}
Zufallsfolge $\cdot$ Lauf $\cdot$ Runs-Test $\cdot$ kumulierte Häufigkeitsverteilung
$\cdot$ exponentielle Verteilung $\cdot$ Zufallszahlengenerator $\cdot$ \texttt{/dev/random}

%-----------------------------------------------------
% 1  Einleitung
%-----------------------------------------------------
\section{Einleitung}\label{sec:einleitung}

Ausgangspunkt dieser Arbeit sind Karl Poppers Untersuchungen zum
Wahrscheinlichkeitsbegriff in \enquote{Logik der Forschung}
\cite{popper1935}: Zufälligkeit ist keine Eigenschaft, die sich an einem
einzelnen Ereignis ablesen lässt, sondern eine Zuschreibung, die an der
Wiederholbarkeit und Überprüfbarkeit von Beobachtungen hängt. Für die hier
behandelten endlichen Bitfolgen ziehen wir daraus die Konsequenz, den
Zufallsbegriff nicht axiomatisch vorauszusetzen, sondern endliche Folgen
anhand beobachtbarer, rein kombinatorischer Strukturen zu untersuchen --
ohne Rückgriff auf Grenzwertaussagen über unendliche Folgen.

In dieser Arbeit wird zuerst ein Auswertungsverfahren für endliche Bitfolgen
entworfen, der Einschritttest (Abschnitt~\ref{sec:einschritttest}), und das
Ergebnis seiner Anwendung auf die Ausgabe eines Zufallsgeräts analysiert
(Abschnitt~\ref{sec:experiment}). Es wird sich experimentell zeigen, dass
die Häufigkeitsverteilung der Läufe und Teilläufe der Ausgabe eines
Zufallsgenerators sich einer exponentiell gearteten Verteilung annähert.
Motiviert durch dieses experimentelle Ergebnis wird die Struktur der
Häufigkeitsverteilungen solcher Bitfolgen formal untersucht und damit die
experimentelle Beobachtung erklärt (Abschnitt~\ref{sec:klassen}). Schließlich
werden Algorithmen angegeben, die Elemente der dabei auftretenden Klassen
erzeugen (Abschnitt~\ref{sec:algorithmen}), und es wird gezeigt, wie sich
diese Klassen über die binäre Darstellung natürlicher Zahlen zusammensetzen
lassen (Abschnitt~\ref{sec:zusammengesetzt}). Abschnitt~\ref{sec:diskussion}
diskutiert die Konsequenzen für den Zufallsbegriff,
Abschnitt~\ref{sec:schluss} fasst zusammen.

Die Beiträge im Einzelnen:
\begin{itemize}
  \item ein lineares Auswertungsverfahren, das die Häufigkeiten von Läufen
        \emph{und} Teilläufen endlicher Bitfolgen kumuliert
        (Einschritttest, Abschnitt~\ref{sec:einschritttest});
  \item eine experimentelle Auswertung der Ausgabe von \texttt{/dev/random}
        (Abschnitt~\ref{sec:experiment});
  \item die Klassen $B_N$ der Bitfolgen mit exakt exponentiell fallender
        kumulierter Häufigkeitsverteilung, inklusive der Struktur ihrer
        Grundverteilung (Satz~\ref{satz:grundverteilung}) und einer
        Abgeltungsidentität für Teilläufe
        (Lemma~\ref{lem:abgeltungsidentitaet});
  \item ein iterativer und ein rekursiver Generator für Elemente aus $B_N$
        (Abschnitt~\ref{sec:algorithmen});
  \item eine eineindeutige Zuordnung zwischen natürlichen Zahlen und
        Zusammensetzungen paarweise verschiedener Klassen $B_N$
        (Abschnitt~\ref{sec:zusammengesetzt}).
\end{itemize}

%-----------------------------------------------------
% 2  Stand der Forschung
%-----------------------------------------------------
\section{Stand der Forschung}\label{sec:stand}

Die mathematische Grundlage des Zufallsbegriffs ist die
Kolmogoroff-Axiomatik \cite{kolmogoroff1933}; der daraus folgende
starke und schwache Grenzwertsatz begründet die Erwartung, dass in
hinreichend langen Folgen unabhängiger, gleichverteilter Bits die
Lauflängen geometrisch verteilt sind: Die Wahrscheinlichkeit einer
Lauflänge $k$ beträgt $2^{-k}$ \cite[und dortige Übungsaufgaben zu Runs]{feller1968}.
Der hier beobachtete exponentielle Abfall der (kumulierten)
Häufigkeiten ist daher für ideale Zufallsfolgen zu erwarten; neu in
dieser Arbeit ist nicht der Abfall selbst, sondern die Kumulation über
die Teilläufe sowie die Untersuchung derjenigen endlichen Folgen, deren
kumulierte Verteilung \emph{exakt} exponentiell fällt.

Statistische Verfahren zur Prüfung von Zufallsfolgen sind seit langem
etabliert. Der klassische Runs-Test von Wald und Wolfowitz
\cite{waldwolfowitz1940} prüft die Anzahl der Läufe einer Folge gegen
ihre erwartete Verteilung; er ist bis heute Bestandteil der
standardisierten Testbatterien für kryptographische Zufallsgeneratoren,
etwa der NIST-Suite \cite{rukhin2010nist} und ihrer Vorläufer
\cite{knuth1997,marsaglia1996}. Übersichten über den Entwurf und die
Prüfung von Zufallszahlengeneratoren geben Knuth \cite{knuth1997} und
L'Ecuyer \cite{lecuyer1994}. Der algorithmische Zufallsbegriff von
Martin-Löf \cite{martinloef1966} definiert Zufälligkeit unendlicher
Folgen über die Unmöglichkeit effektiver Tests; unserem Zugang am
nächsten steht die dortige Idee, Zufälligkeit über das Bestehen
effektiver Tests zu fassen, mit dem Unterschied, dass wir endliche
Folgen und eine einzige, rein kombinatorische Kenngröße betrachten.

Als experimentelle Quelle dient das Gerät \texttt{/dev/random} einer
Linux-Installation; seine Arbeitsweise ist dokumentiert und mehrfach
analysiert worden \cite{gutterman2004}. Die hier verwendete
Auswertung ist unabhängig von der speziellen Quelle; wir greifen
\texttt{/dev/random} nur als notorisch guten Lieferanten für Bitfolgen
heraus, deren Zustandekommen nicht kausal erklärt werden soll.

Die erkenntnistheoretische Rahmung folgt Popper \cite{popper1935}.
Auf die Frage, ob und in welchem Sinne die Klassen $B_N$ den
wahrscheinlichkeitstheoretischen Zufallsbegriff ersetzen können, kommen
wir in Abschnitt~\ref{sec:diskussion} zurück.

%-----------------------------------------------------
% 3  Grundlagen: Läufe und Teilläufe
%-----------------------------------------------------
\section{Grundlagen: Läufe und Teilläufe}\label{sec:grundlagen}

\begin{defi}[Menge der endlichen Bitfolgen]\label{defi:bitfolgen}
\begin{equation}\label{eqn:mengederbitfolgen}
  \mathcal B := \left\{(b_m)_{1\leq m\leq M<\infty},\ b_m\in\left\{0,1\right\}\right\}
\end{equation}
\end{defi}

\noindent Wir gehen von einer beliebigen endlichen Bitfolge
\begin{equation}\label{eqn:beliebigebitfolge}
  (b_m)_{1\leq m\leq M,\ M<\infty}\in\mathcal B
\end{equation}
aus. Zum Beispiel
\begin{equation}\label{eqn:runtestfolge}
  10001010100111110000011111110001010101
\end{equation}
Diese Bitfolge können wir auch als Folge von Läufen auffassen. Dabei ist
ein Lauf eine Folge von gleichbleibenden Bits. In unserem Beispiel sind
es
\begin{equation}\label{eqn:runtestfolgevonlaeufen}
  1\cdot 000\cdot 1\cdot 0\cdot 1\cdot 0\cdot 1\cdot 00\cdot 11111\cdot
  00000\cdot 1111111\cdot 000\cdot 1\cdot 0\cdot 1\cdot 0\cdot 1\cdot 0\cdot 1
\end{equation}
\begin{defi}[Lauf]\label{defi:laeufe}
Ein \emph{Lauf} einer Bitfolge $(b_m)_{1\leq m\leq M}\in\mathcal B$ ist eine
konstante, in beiden Richtungen maximale Teilfolge $(b_p,\dots,b_q)$ mit
$1\leq p\leq q\leq M$:
\begin{equation}\label{eqn:mengederlaeufe}
  b_p=b_{p+1}=\dots=b_q,\qquad
  \left(p=1\ \text{oder}\ b_{p-1}\neq b_p\right),\qquad
  \left(q=M\ \text{oder}\ b_q\neq b_{q+1}\right).
\end{equation}
Die Zahl $q-p+1$ heißt die \emph{Länge} des Laufs.
\end{defi}

\begin{lemma}[Läufzerlegung]\label{lem:lzerlegung}
Jede Bitfolge $(b_m)_{1\leq m\leq M}\in\mathcal B$ zerfällt eindeutig in
Läufe. Mit $m_1:=1$ und der rekursiven Vorschrift
\begin{equation}\label{eqn:zerlegung}
  m_{i+1}:=\min\left(\left\{m>m_i:\ b_m\neq b_{m_i}\right\}\cup\{M+1\}\right)
\end{equation}
ergeben sich eindeutige Zerlegungspunkte
$1=m_1<m_2<\dots<m_N<m_{N+1}=M+1$ mit Läufen
$(b_{m_i},\dots,b_{m_{i+1}-1})$ und Laufängen $r_n:=m_{n+1}-m_n$.
Insbesondere gilt
\begin{equation}
  \sum_{n=1}^{N} r_n=M,
\end{equation}
und aufeinanderfolgende Läufe tragen komplementäre Bitwerte. Die Folge
$(r_n)_{1\leq n\leq N}\in\mathbb N$ heißt die \emph{Lauffolge} von
$(b_m)$.
\end{lemma}

\begin{proof}
Da $(b_m)$ auf jedem Block $[m_i,m_{i+1})$ konstant ist und
$b_{m_{i+1}}\neq b_{m_i}$ gilt (sofern $m_{i+1}\leq M$), ist
(\ref{eqn:zerlegung}) wohldefiniert und liefert eine streng wachsende
Folge, die nach endlich vielen Schritten $M+1$ erreicht. Die
Maximalitätsbedingungen aus (\ref{eqn:mengederlaeufe}) sichern, dass
jeder Block ein Lauf im Sinne der Definition ist. Die Eindeutigkeit
folgt daraus, dass der erste Lauf bei $m=1$ beginnen muss und jeder
Folgelauf durch das Ende des vorherigen festgelegt ist.
\end{proof}

\begin{bem}[Komplement und Lauffolge]\label{bem:komplement}
Zwei Bitfolgen besitzen genau dann dieselbe Lauffolge, wenn sie gleich
oder binäre Komplemente voneinander sind: Die Lauffolge legt -- zusammen
mit dem Bit des ersten Laufs -- wegen des alternierenden Aufbaus die
Bitfolge eindeutig fest, und der Übergang zum Komplement ändert nur
dieses eine Bit.
\end{bem}

\noindent Da in sämtlichen Häufigkeitsbetrachtungen dieser Arbeit nur die
Längen der Läufe eingehen, identifizieren wir Läufe fortan mit ihrer
Länge und arbeiten mit Lauffolgen $(r_n)\in\mathbb N$.

Des Weiteren führen wir die triviale Überlegung an, dass ein Lauf der
Länge $N$ auch Läufe $s<N$ mit entsprechender Häufigkeitsverteilung
\begin{equation}\label{eqn:haufigkeitenteillaeufe}
  _NS=(_Ns_n)_{0<n\leq N},\quad\text{mit}\ _Ns_n=N-n+1
\end{equation}
beinhaltet. Wir nehmen zur Illustration den Lauf $N=8$
\begin{equation}\label{eqn:teillaufindemlauf}
  1\;1\;1\;1\;1\;1\;1\;1\quad\text{bzw.}\quad 0\;0\;0\;0\;0\;0\;0\;0
\end{equation}
Dann beinhaltet $N$:
\begin{align*}
  {}_8s_1&=8-1+1=8\ \text{Läufe } 1, & {}_8s_2&=8-2+1=7\ \text{Läufe } 2,\\
  {}_8s_3&=8-3+1=6\ \text{Läufe } 3, & {}_8s_4&=8-4+1=5\ \text{Läufe } 4,\\
  {}_8s_5&=8-5+1=4\ \text{Läufe } 5, & {}_8s_6&=8-6+1=3\ \text{Läufe } 6,\\
  {}_8s_7&=8-7+1=2\ \text{Läufe } 7, & {}_8s_8&=8-8+1=1\ \text{Lauf } 8.
\end{align*}

\begin{defi}[Teillauf]\label{defi:teillauf}
Die in einem Lauf enthaltenen Läufe kleinerer Länge heißen fortan
Teilläufe.
\end{defi}

%-----------------------------------------------------
% 4  Der Einschritttest
%-----------------------------------------------------
\section{Der Einschritttest}\label{sec:einschritttest}

Wir fassen die zu untersuchende Bitfolge als Folge von Läufen
\begin{equation}\label{eqn:laeuflangefolge}
  (r_n)_{0<n\leq N},\quad r_n\in\mathcal \mathbb N,\quad N\text{ geeignet}
\end{equation}
auf. Um eine solche Folge auszuwerten, ermitteln wir die kumulierte
Verteilung
\begin{equation}\label{eqn:verteilunglaeufeteillaeufe}
  L=(l_p)_{0<p\leq P<\infty}
\end{equation}
der Häufigkeiten $l_p$ der Läufe und der darin enthaltenen Teilläufe der
Folge $(r_n)_{0<n\leq N<\infty}$.

Motivieren wir diesen Ansatz durch die folgende Überlegung: Für die
Ausgabe $(b_m)\in \mathcal B$ eines Zufallsalgorithmus, der unabhängige,
gleichverteilte Bits liefert, gilt
$P(b_{n-1}=b_n)=P(b_{n-1}\neq b_n)=\tfrac{1}{2}$ \cite{feller1968}.
Statt auf Grenzwertaussagen zurückzugreifen, bevorzugen wir es, für
jedes Bit $b_m$ aus $(b_m)$ die Häufigkeit der Fälle
$(b_{m})=(b_{m+i})$ zu betrachten. Dies führt uns zu dem Ansatz der
kumulierten Verteilungen: Wir fassen $(b_m)\in\mathcal B$ als Folge
$(r_n)$ auf und ermitteln die Häufigkeiten ihrer Läufe und
der darin enthaltenen Teilläufe.

Es gilt auf triviale Weise, dass der größte Index $P$ der Folge
(\ref{eqn:verteilunglaeufeteillaeufe}) durch den längsten Lauf von
$(r_n)$ gegeben ist:
\begin{equation}\label{eqn:maximalelaengevonl}
  P=\max\left\{ r_1,\dots, r_N\right\}
\end{equation}

\noindent Aus den vorangegangenen Überlegungen ergibt sich die folgende
Vorschrift, die wir hier mit der Bezeichnung Einschritttest versehen.

\begin{figure}[htbp]
\begin{lstlisting}[caption={Einschritttest für Folgen von Läufen},label={algo:einschritt-laeufe}]
Eingabe: Eine beliebige Folge von Läufen der Länge N, N endlich
Lege einen Array L der Länge maximal N an
fuer jeden Lauf der Folge
    ermittle die Laenge des Laufs (hier l)
    fuer jede Stelle i = 1 bis l
        erhoehe die Haeufigkeit an der Stelle i um l - i + 1
Gib L aus
\end{lstlisting}
\end{figure}

\noindent Man kann dasselbe Verfahren auf die der Folge von Läufen
$(r_n)$ zugrunde liegende Bitfolge $(b_m)$ bzw.
$(\overline b_m)$ anwenden. Es ergibt sich eine etwas andere
Rechenvorschrift. Sie liefert dasselbe Ergebnis wie
Algorithmus~\ref{algo:einschritt-laeufe}.
% TODO(Inhalt): Pseudocode präzisieren: Die innere Schleife muss fuer JEDES
% Bit (nicht pro Lauf) durchlaufen werden, sonst zaehlt der Algorithmus
% Laenge>=i statt Teillaeufe und ist Algorithmus 1 nicht aequivalent.
% Beide Listings durch getesteten Referenzcode (z. B. C/Python) ersetzen und
% Aequivalenz beweisen oder als Lemma zeigen.
\begin{figure}[htbp]
\begin{lstlisting}[caption={Einschritttest für Bitfolgen},label={algo:einschritt-bitfolge}]
Eingabe: Eine beliebige endliche Folge von Bits der Laenge P, P endlich
Lege einen Array L der Laenge maximal P an
Lege einen Zaehler z mit dem Wert 0 an
fuer jeden Bit der Folge
    hat sich der Wert des Bits im Vergleich zum vorherigen Schritt geaendert?
    wenn nein, dann erhoehe den Wert des Zaehlers z um 1
    wenn ja, dann stelle den Zaehler z auf 1 zurueck
    fuer jede Stelle i = 1 bis z
        erhoehe die Haeufigkeit an der Stelle i um 1
Gib L aus
\end{lstlisting}
\end{figure}
Der Name \enquote{Einschritttest} wurde aus der Tatsache hergeleitet,
dass in der Auffassung des Algorithmus~\ref{algo:einschritt-bitfolge},
welche chronologisch als erste entstanden ist, die Eingabe Schritt für
Schritt abgearbeitet wird.

%-----------------------------------------------------
% 5  Experiment
%-----------------------------------------------------
\section{Ein motivierendes Experiment}\label{sec:experiment}

Wir greifen auf ein experimentelles Mittel zurück, um eine Bitfolge zu
erzeugen, deren Struktur nicht kausal erklärt werden soll: eine
sogenannte Zufallsfolge. Ein solches experimentelles Mittel ist das
Gerät \texttt{/dev/random} einer beliebigen Linux-Installation
\cite{gutterman2004}, das den Ruf besitzt, ein notorisch guter
Lieferant solcher Bitfolgen zu sein, auch wenn es relativ langsam
arbeitet. Darüber hinaus kann seine Arbeitsweise lückenlos nachgeprüft
werden. Wir sammeln die Ausgabe von \texttt{/dev/random} auf einer oder
mehreren Linux-Instanzen bis zu einem Volumen von ca.\
$16\cdot 10^6$ Bits.
% TODO(Inhalt): Exakte Bitanzahl M angeben (nicht nur "ca."); nach
% Definition des Einschritttests muss l_1 = M exakt gelten.
% Zudem: Die Zeilen fuer die Lngen 8--11 der Tabelle brechen die
% Halbierungsreihe offensichtlich (vgl. erwartete Werte ~126000, ~63000,
% ~31000, ~16000). Daten neu erheben oder Erklrung angeben; Erfassungsskript
% und Rohdaten veroeffentlichen.
Das Ergebnis dieses experimentellen Vorgehens ist eine endliche
Bitfolge, deren Länge $M$ durch das Experiment bedingt ist:
\begin{equation}\label{eqn:bitfolgebeliebigerlaenge}
  (b_m)_{0<m\leq M<\infty},\quad b_m\in\{0,1\}
\end{equation}

\noindent Wie in Abschnitt~\ref{sec:grundlagen} über Läufe und
Teilläufe ausgeführt, fassen wir nun die Folge $(b_m)$ aus
(\ref{eqn:bitfolgebeliebigerlaenge}) als eine Folge gemäß
(\ref{eqn:laeuflangefolge}) auf und werten sie mit dem Einschritttest
aus. Das Ergebnis der Auswertung der Ausgabe des Geräts
\texttt{/dev/random} ist in Tabelle~\ref{tab:haeufigkeiten}
protokolliert. Es fällt sofort auf, dass die Verteilung der Läufe und
Teilläufe der Ausgabe des Geräts \texttt{/dev/random} exponentiell
fallend verläuft. Wir stellen diese Ergebnisse in
Abbildung~\ref{fig:haeufigkeiten} graphisch dar. Die dort und in
Tabelle~\ref{tab:haeufigkeiten} als theoretische Werte ausgewiesene
Reihe besteht aus den Potenzen der Funktion $2^{25-i}$ und dient der
Veranschaulichung eines exponentiell abfallenden Verlaufs. Im Sinne der
intersubjektiven Nachprüfung dieses Experimentes wird der Leser
ausdrücklich dazu ermutigt, dieses Experiment mit den Mitteln seiner
Wahl nachzustellen. Wie in unserem Fall sollte der Ausgang des
Experimentes einen tendenziell exponentiell abfallenden Verlauf
aufweisen.

\begin{table}[htbp]
\centering
\begin{tabular}{rrrrrrrr}
\toprule
Länge & Häufigkeit & $\log_2$ & theo.\ Wert & Länge & Häufigkeit & $\log_2$ & theo.\ Wert\\
\midrule
$1$&$16170866$&$23{,}95$&$24$&$13$&$4021$&$11{,}97$&$12$\\
$2$&$8087159$&$22{,}95$&$23$&$14$&$2034$&$10{,}99$&$11$\\
$3$&$4044093$&$21{,}95$&$22$&$15$&$1011$&$9{,}98$&$10$\\
$4$&$2021615$&$20{,}95$&$21$&$16$&$504$&$8{,}98$&$9$\\
$5$&$1010537$&$19{,}95$&$20$&$17$&$251$&$7{,}97$&$8$\\
$6$&$494406$&$18{,}95$&$19$&$18$&$133$&$7{,}06$&$7$\\
$7$&$251941$&$17{,}94$&$18$&$19$&$69$&$6{,}11$&$6$\\
$8$&$25921$&$16{,}94$&$17$&$20$&$33$&$5{,}04$&$5$\\
$9$&$2953$&$15{,}94$&$16$&$21$&$15$&$3{,}91$&$4$\\
$10$&$1469$&$14{,}94$&$15$&$22$&$7$&$2{,}81$&$3$\\
$11$&$5776$&$13{,}95$&$14$&&&&\\
$12$&$7969$&$12{,}96$&$13$&&&&\\
\bottomrule
\end{tabular}
\caption{Häufigkeit der Läufe und Teilläufe der Ausgabe von
\texttt{/dev/random} einer beliebigen Linux-Installation. In der
zweiten Spalte werden die tatsächlichen Häufigkeiten aufgelistet, in
der dritten ihre Entsprechung auf einer logarithmischen Skala zur
Basis $2$. Die theoretischen Werte sind jene der Funktion $2^{25-n}$.}
\label{tab:haeufigkeiten}
\end{table}

\begin{figure}[htbp]
\centering
\IfFileExists{experiment.png}{%
  \includegraphics[width=1.0\textwidth]{experiment.png}%
}{%
  \fbox{\parbox[c][6cm][c]{0.9\textwidth}{\centering
    Die Abbildungsdatei \texttt{experiment.png} fehlt.\\[0.5em]
    \emph{TODO(Inhalt):} Graphik aus den Daten der
    Tabelle~\ref{tab:haeufigkeiten} erzeugen (Vektorgrafik, z.\,B.\
    pgfplots; doppelt-logarithmische Darstellung mit Fehlerbalken).}}%
}
\caption{Graphische Darstellung der Ergebnisse aus
Tabelle~\ref{tab:haeufigkeiten}.}
\label{fig:haeufigkeiten}
\end{figure}

Diese Beobachtung lässt vermuten, dass die Ausgabe des Geräts
\texttt{/dev/random} -- oder irgendeines Geräts, das den Zweck erfüllt,
Bitfolgen ohne kausale Erzeugung zu liefern -- einen annähernd
exponentiell fallenden Verlauf der Verteilung ihrer Läufe und
enthaltenen Teilläufe aufweist. Widersprüchlicherweise wäre damit doch
eine Struktur des Zufalls erklärt. Wir widmen uns in den nächsten
Abschnitten der Untersuchung solcher Bitfolgen, ohne Rücksicht auf die
Art, wie sie gewonnen werden.

%-----------------------------------------------------
% 6  Bitfolgen mit exponentiell kumulierten Haeufigkeitsverteilungen
%-----------------------------------------------------
\section{Bitfolgen mit exponentiell kumulierten
Häufigkeitsverteilungen}\label{sec:klassen}

\subsection{Kumulierte Verteilung und die Klassen \texorpdfstring{$B_N$}{B\_N}}

\begin{defi}[Kumulierte Verteilung]\label{defi:kumulierteverteilung}
Sei $(_Nb_m)$ eine Bitfolge aus $\mathcal B$ und
$(r_n)_{0<n\leq N<\infty}$ ihre entsprechende Folge von Läufen. Wir
bezeichnen mit der Folge
\begin{equation}\label{eqn:defkumuliert}
  _NL:=(\,_Nl_i\,)_{0<i\leq N<\infty}
\end{equation}
die Verteilung der Häufigkeiten der Läufe und enthaltenen Teilläufe
der Bitfolge $(_Nb_m)$ und nennen sie die kumulierte Verteilung.
\end{defi}

\begin{defi}[Exponentiell fallende kumulierte Verteilung]\label{defi:exponentiell}
Unter der exponentiell fallenden, kumulierten Verteilung
$_NL:=(\,l_i\,)_{0<i\leq N<\infty}$ einer Bitfolge
$(_Nb_m)\in\mathcal B$ verstehen wir
\begin{equation}\label{eqn:exponentiellerverlauf}
  \forall\ 0<n\leq N<\infty:\quad _Nl_n=2^{N-n}
\end{equation}
Äquivalent dazu ist
\begin{align}
  _Nl_N&=1\label{eqn:ln gleich 1}\\
  _Nl_{n-1}&=2\ {_Nl}_n,\quad 2\leq n\leq N\label{eqn:lnminus1 gleich 2 ln}
\end{align}
\end{defi}

\noindent Wir machen darauf aufmerksam, dass die kumulierte Verteilung
einer Bitfolge nicht hinreicht, um diese Folge vollständig zu
charakterisieren. Vielmehr entspricht diese Verteilung einer Klasse von
Bitfolgen. Daher die folgende Definition.

\begin{defi}[Klassen $B_N$]\label{defi:bn}
Wir bezeichnen mit $B_N$ die Menge der Folgen mit längstem Lauf $N$ und
mit exponentiell fallender, kumulierter Verteilung gemäß
(\ref{eqn:exponentiellerverlauf}) bzw.~(\ref{eqn:ln gleich 1}) und
(\ref{eqn:lnminus1 gleich 2 ln}).
\end{defi}

\begin{defi}[Menge der Klassen]\label{defi:bz}
Menge der Klassen der Bitfolgen mit exponentiell fallenden, kumulierten
Verteilungen:
\begin{equation}\label{eqn:mengederklassen}
  \mathcal{B_Z} := \{B_N,\ 0< N<\infty \}
\end{equation}
\end{defi}

\subsection{Die Grundverteilung}

Nachdem die grundlegenden Definitionen gestellt wurden, können wir uns
der Untersuchung der Elemente aus einem beliebigen $B_N$ widmen. Aus
der Definition der exponentiell fallenden, kumulierten Verteilung
ausgehend kann untersucht werden, ob die Verteilung der Häufigkeiten
ihrer Läufe ohne Teilläufe einer Systematik unterliegt. Wir rechnen
also aus jeder Häufigkeit der kumulierten Verteilung $_NL$ ihre bereits
abgegoltenen Teilläufe heraus; das Ergebnis nennen wir die
Grundverteilung.

\begin{defi}[Grundverteilung]\label{defi:grundverteilung}
Sei $(b_m)\in\mathcal B$ eine Bitfolge mit Lauffolge $(r_n)_{1\leq n\leq R}$
und kumulierter Verteilung $L=(l_p)_{1\leq p\leq P}$. Die
\emph{Grundverteilung} $G=(g_p)_{1\leq p\leq P}$ zählt die Läufe je Länge,
\begin{equation}\label{eqn:grunddef}
  g_p:=\left|\left\{n\in\{1,\dots,R\}:\ r_n=p\right\}\right|,\qquad 1\leq p\leq P.
\end{equation}
Kumulierte und Grundverteilung hängen über die Dreiecksbeziehung
\begin{equation}\label{eqn:kumuliertrelation}
  l_p=\sum_{k=p}^{P}\left(k-p+1\right)\cdot g_k,\qquad 1\leq p\leq P,
\end{equation}
zusammen: Jeder Lauf der Länge $k$ enthält genau $k-p+1$ Teilläufe der
Länge $p$. Wegen des Dreiecksaufbaus von (\ref{eqn:kumuliertrelation}) ist
$G$ durch $L$ eindeutig bestimmt; die rekursive Auflösung von oben,
\begin{equation}\label{eqn:abgeltungrekursiv}
  g_P=l_P,\qquad
  g_p=l_p-\sum_{k=p+1}^{P}\left(k-p+1\right)\cdot g_k
  \quad(p=P-1,\dots,1),
\end{equation}
ist genau das im Folgenden beschriebene Abgeltungsverfahren.
\end{defi}

Zur Veranschaulichung vollziehen wir die Auflösung
(\ref{eqn:abgeltungrekursiv}) an einem Beispiel aus $B_5$ nach. In diesem
Fall betragen die Werte von $_5L=(_5l_i)_{0<i\leq5}$
\begin{equation}
  _5l_i=2^{5-i}
\end{equation}
Das heißt explizit, dass
\begin{itemize}
  \item $_5l_1=16$,
  \item $_5l_2=8$,
  \item $_5l_3=4$,
  \item $_5l_4=2$,
  \item $_5l_5=1$.
\end{itemize}
Der längste Lauf beträgt $5$ und ist kein Teillauf. Wir ziehen also die
durch den Lauf $_5l_5$ abgegoltenen Teilläufe ab und erhalten:
\begin{itemize}
  \item $5$ Teilläufe $1$,
  \item $4$ Teilläufe $2$,
  \item $3$ Teilläufe $3$,
  \item $2$ Teilläufe $4$,
  \item $1$ Lauf $5$.
\end{itemize}
Nun ist der Lauf $5$ abgegolten und damit auch beide Läufe $4$. Noch
geltend zu machen sind die folgenden Läufe:
\begin{itemize}
  \item $_5l_1-5=11$ Läufe $1$,
  \item $_5l_2-4=4$ Läufe $2$,
  \item $_5l_3-3=1$ Läufe $3$,
  \item $_5l_4-2=0$ Läufe $4$,
  \item $_5l_5-1=0$ Läufe $5$.
\end{itemize}
Wir machen nun den Lauf $3$ geltend, da er nun nicht mehr als Teillauf
vorhanden ist. Das heißt
\begin{itemize}
  \item $3$ Teilläufe $1$,
  \item $2$ Teilläufe $2$,
  \item $1$ Lauf $3$.
\end{itemize}
Nun ist der Lauf $3$ abgegolten. Noch geltend zu machen sind die
folgenden Läufe:
\begin{itemize}
  \item $_5l_1-5-3=8$ Läufe $1$,
  \item $_5l_2-4-2=2$ Läufe $2$,
  \item $_5l_3-3-1=0$ Läufe $3$.
\end{itemize}
Wir machen nun die Läufe $2$ geltend. Das heißt
\begin{itemize}
  \item $4$ Teilläufe $1$,
  \item $2$ Teilläufe $2$.
\end{itemize}
Nun ist der Lauf $2$ abgegolten. Noch geltend zu machen sind die
folgenden Läufe:
\begin{itemize}
  \item $_5l_1-5-3-4=4$ Läufe $1$.
\end{itemize}
Da die Läufe $1$ keine Teilläufe enthalten, werden sie unmittelbar
abgegolten. Das Ergebnis der Auflösung (\ref{eqn:abgeltungrekursiv}) ist
die Grundverteilung
\begin{equation}
  _5G=(\;4;2;1;0;1\;).
\end{equation}
Eine mögliche Realisierung dieser Verteilung ist
\begin{equation}\label{eqn:folfgeausb5}
  (\;0101001100011111\;)\ \text{mit}\
  (\;\underbrace{ 0\; 1\; 0\; 1}_{4}\;
   \underbrace{00\; 11}_{2}\;
   \underbrace{000}_{1}\;
   \underbrace{\;\;\;\;}_{0}\;
   \underbrace{11111}_{1}\; )\in B_5
\end{equation}
Natürlich ist jede andere Permutation dieser Folge von Läufen gültig.
Dabei werden bei Bedarf die Werte der Bits komplementiert. Wir
wiederholen beispielhaft diese Berechnung für Bitfolgen aus den Klassen
$B_1$ bis $B_8$:
\begin{align}
  (\;1\;)_{[B_1]}\label{eqn:folfgeausb1}\\
  (\;0\;;\;1\;)_{[B_2]}\nonumber\\
  (\;1\;;\;0\;;\;1\;)_{[B_3]}\nonumber\\
  (\;2\;;\;1\;;\;0\;;\;1\;)_{[B_4]}\nonumber\\
  (\;4\;;\;2\;;\;1\;;\;0\;;\;1\;)_{[B_5]}\nonumber\\
  (\;8\;;\;4\;;\;2\;;\;1\;;\;0\;;\;1\;)_{[B_6]}\nonumber\\
  (\;16\;;\;8\;;\;4\;;\;2\;;\;1\;;\;0\;;\;1\;)_{[B_7]}\nonumber\\
  (\;32\;;\;16\;;\;8\;;\;4\;;\;2\;;\;1\;;\;0\;;\;1\;)_{[B_8]}
  \label{eqn:folfgeausb8}
\end{align}
Wir bemerken sofort, dass mit wachsendem $N$ die Grundverteilung der
Elemente aus den $B_N$ eine gewisse Struktur aufweist: Lediglich eine
Verschiebung nach rechts der Häufigkeit der Läufe $2$ bis $N$ und eine
Verdoppelung der Häufigkeit der Läufe $1$ sind zu verzeichnen. Wir wollen
nun diese experimentellen Beobachtungen formal untermauern.

\begin{satz}\label{satz:grundverteilung}
Sei $N\geq 2$ und $(_Nb_m)\in B_N$. Dann gilt für die Grundverteilung
$(_Ng_n)_{1\leq n\leq N}$ der Läufe von $(_Nb_m)$
\begin{align}
  _Ng_n&=2^{N-2-n},\quad 1\leq n\leq N-2,\label{eqn:gnform}\\
  _Ng_{N-1}&=0,\label{eqn:g_nminus1}\\
  _Ng_N&=1.\label{eqn:g_n}
\end{align}
Insbesondere erfüllt $_NG$ die Rekursion $_Ng_n=2\,{_Ng}_{n+1}$ für
$1\leq n\leq N-3$. (Für $N=1$ gilt $_1g_1=1$.)
\end{satz}

\begin{proof}
Da der längste Lauf die Länge $N$ hat, ist $P=N$, und es ist
$_Nl_n=2^{N-n}$ für alle $1\leq n\leq N$
(Gleichung (\ref{eqn:exponentiellerverlauf})). Wir lösen die Rekursion
(\ref{eqn:abgeltungrekursiv}) von oben nach unten auf.

Für $n=N$ ist $_Ng_N=_Nl_N=2^{0}=1$, das ist (\ref{eqn:g_n}).
Für $n=N-1$ ist
\begin{equation*}
  _Ng_{N-1}=_Nl_{N-1}-2\cdot _Ng_N=2-2=0,
\end{equation*}
das ist (\ref{eqn:g_nminus1}).

Für die verbleibenden Indizes zeigen wir die geschlossene Form
(\ref{eqn:gnform}) durch Abwärtsinduktion über $n=N-2,\dots,1$.
Subtraktion der Gleichungen (\ref{eqn:kumuliertrelation}) für $n$ und
$n+1$ liefert dabei
\begin{equation}\label{eqn:differenz}
  _Nl_n-\ _Nl_{n+1}= _Ng_n+\sum_{k=n+1}^{N} _Ng_k,\qquad 1\leq n\leq N-1,
\end{equation}
denn der Beitrag eines Laufs der Länge $k>n$ an $l_n$ übersteigt seinen
Beitrag an $l_{n+1}$ um genau $1$.

\emph{Induktionsanfang} $n=N-2$: Aus (\ref{eqn:abgeltungrekursiv})
folgt
\begin{equation*}
  _Ng_{N-2}=_Nl_{N-2}-2\cdot\underbrace{_Ng_{N-1}}_{=0}
  -3\cdot\underbrace{_Ng_N}_{=1}
  =4-3=1=2^{N-2-(N-2)}.
\end{equation*}

\emph{Induktionsschritt}: Sei $n\leq N-3$ und es gelte bereits
$_Ng_k=2^{N-2-k}$ für alle $n<k\leq N-2$. Aus (\ref{eqn:differenz}) und
$\,__Nl_n-\ _Nl_{n+1}=2^{N-n-1}$ folgt dann mit der geometrischen Summe
\begin{equation*}
  _Ng_n
  =2^{N-n-1}-\sum_{k=n+1}^{N-2}2^{N-2-k}
  -\underbrace{_Ng_{N-1}}_{=0}-\underbrace{_Ng_N}_{=1}
  =2^{N-n-1}-\left(2^{N-2-n}-1\right)-1
  =2^{N-2-n}.
\end{equation*}
Das zeigt (\ref{eqn:gnform}) für alle $1\leq n\leq N-2$ und damit die
Behauptung.
\end{proof}

\begin{bem}[Besonderheit des Verfahrens]\label{bem:keinlaufnminus1}
In keinem Element aus $B_N$, $N\geq 2$, tritt ein Lauf der Länge $N-1$
auf: $_Ng_{N-1}=0$. Die Grundverteilung springt unmittelbar von
$_Ng_{N-2}=1$ auf $_Ng_N=1$. Das ist keine Unregelmäßigkeit, sondern
eine besondere, systematische Eigenschaft des Verfahrens: Der
Häufigkeitsanteil der Länge $N-1$ ist in der Kumulation vollständig in
den Teilläufen des einzigen Laufs der Länge $N$ aufgegangen --
$_Nl_{N-1}=2$ zählt genau dessen beide Teilläufe der Länge $N-1$ --,
und bei der Auflösung von oben wird dieser Anteil abgezogen, bevor ein
eigener Lauf der Länge $N-1$ frei werden könnte.
\end{bem}

\begin{lemma}[Summenidentität]\label{lem:summenidentitaet}
Für alle $m\geq 0$ gilt
\begin{equation}\label{eqn:summenidentitaet}
  T_m:=\sum_{k=1}^{m} k\cdot 2^{m-k}=2^{m+1}-m-2.
\end{equation}
\end{lemma}

\begin{proof}
Induktion über $m$. Für $m=0$ ist die Summe leer und
$T_0=0=2^{1}-0-2$. Im Schritt gilt
\begin{equation*}
  T_{m+1}=\sum_{k=1}^{m+1}k\cdot 2^{m+1-k}
  =2\sum_{k=1}^{m}k\cdot 2^{m-k}+(m+1)
  =2T_m+(m+1)
  =2^{m+2}-(m+1)-2.\qedhere
\end{equation*}
\end{proof}

\begin{korollar}[Umkehrung und Existenz]\label{kor:umkehrung}
Besitzt eine Bitfolge die Laufverteilung
\begin{equation}\label{eqn:grundform}
  _Ng_n=2^{N-2-n}\ (1\leq n\leq N-2),\qquad _Ng_{N-1}=0,\qquad _Ng_N=1,
\end{equation}
so liegt sie in $B_N$. Insbesondere ist $B_N$ für jedes $N\geq 1$
nichtleer.
\end{korollar}

\begin{proof}
Zu zeigen ist, dass die kumulierte Verteilung $_Nl_n=2^{N-n}$ beträgt.
Aus (\ref{eqn:kumuliertrelation}) und (\ref{eqn:grundform}) folgt für
$1\leq n\leq N-2$ mit $m:=N-2-n$
\begin{align*}
  _Nl_n&=\sum_{k=n}^{N-2}\left(k-n+1\right)\cdot 2^{N-2-k}
  +\left(N-n+1\right)\cdot 1\\
  &=\sum_{j=0}^{m}\left(j+1\right)\cdot 2^{m-j}+\left(m+3\right)
  =\sum_{k=1}^{m+1} k\cdot 2^{m+1-k}+\left(m+3\right)\\
  &=\left(2^{m+2}-m-3\right)+\left(m+3\right)=2^{m+2}=2^{N-n},
\end{align*}
wobei der vorletzte Schritt die Summenidentität
(Lemma~\ref{lem:summenidentitaet}, $m\mapsto m+1$) benutzt. Für
$n=N-1$ ist $_Nl_{N-1}=2\cdot _Ng_N=2=2^{N-(N-1)}$ und für $n=N$ ist
$_Nl_N=_Ng_N=1$. Der längste Lauf hat die Länge $N$, also liegt die
Bitfolge in $B_N$.
\end{proof}

\subsection{Die Abgeltung der Teilläufe}

Die Abgeltung von Teilläufen, wie sie in dem Verfahren zur Herleitung
der Grundverteilung dargestellt wurde, kann auch direkt berechnet
werden: Die Häufigkeit $_Nt_n$ der abgegoltenen Teilläufe der Länge $n$
ist die Summe der Teillauf-Beiträge aller längeren Läufe; ein Lauf der
Länge $i$ enthält genau $i-n+1$ Teilläufe der Länge $n$:
\begin{equation}\label{eqn:abgeltung}
  _Nt_n=\sum_{i=n+1}^{N}\left(i-n+1\right)\cdot _Ng_i,\qquad 1\leq n\leq N.
\end{equation}
Mit $_Nt_n = _Nl_n - _Ng_n$ ergibt sich der allgemeine Ausdruck
\begin{equation}\label{eqn:abgegoltete}
  _Nl_n - _Ng_n=\sum_{i=n+1}^{N}\left(i-n+1\right)\cdot _Ng_i,
\end{equation}
der für beliebige Grundverteilungen gilt und die Dreiecksbeziehung
(\ref{eqn:kumuliertrelation}) äquivalent umformt. Für die exponentiell
abfallenden kumulierten Verteilungen $_Nl_n=2^{N-n}$ und die
Grundverteilung aus Satz~\ref{satz:grundverteilung} rechnen wir drei
Fälle nach.

\noindent$\blacklozenge$ $n=N$: Die Summe in (\ref{eqn:abgeltung}) ist
leer, also
\begin{equation}
  _Nt_N=0.
\end{equation}

\noindent$\blacklozenge$ $n=N-1$:
\begin{equation}
  _Nt_{N-1}=\left(N-(N-1)+1\right)\cdot\underbrace{_Ng_N}_{=1}=2
  =2^{N-(N-1)}-\underbrace{_Ng_{N-1}}_{=0}.
\end{equation}

\noindent$\blacklozenge$ $1\leq n\leq N-2$: Mit $_Nl_n=2^{N-n}$ und
$_Ng_n=2^{N-2-n}$ gilt
\begin{equation}\label{eqn:tnformel}
  _Nt_n=_Nl_n-_Ng_n=2^{N-n}-2^{N-2-n}=3\cdot 2^{N-2-n}.
\end{equation}
Aus (\ref{eqn:tnformel}) ergibt sich nach Indexverschiebung die
Verschiebungsstruktur, die auch in den Beispielen
(\ref{eqn:folfgeausb1}) bis (\ref{eqn:folfgeausb8}) sichtbar ist: Für
$1\leq n\leq N-2$ gilt
\begin{equation}\label{eqn:verschiebung}
  _Nt_n=3\cdot 2^{N-2-n}=\,_{N-n+1}t_1,
\end{equation}
denn nach (\ref{eqn:tnformel}) ist $_Mt_1=3\cdot 2^{M-3}$, also für
$M=N-n+1$ gerade die rechte Seite von (\ref{eqn:verschiebung}). Die
abgegoltenen Teilläufe der Länge $n$ einer Klasse $B_N$ stimmen also
mit den abgegoltenen Teilläufen der Länge $1$ der Klasse $B_{N-n+1}$
überein.

Zusammenfassend notieren wir:

\begin{lemma}\label{lem:abgeltungsidentitaet}
Sei $N<\infty$. Dann gilt
\begin{equation}\label{eqn:abgeltungsidentitaet}
  2^{N} = \left(N+1\right)+\left( 1_{\{N>1\}}\cdot\sum_{n=0}^{N-2} (N-1-n) \cdot 2^n\right)
\end{equation}
\end{lemma}

\begin{proof}
Die Substitution $k=N-1-n$ führt die Summe über in
\begin{equation*}
  \sum_{n=0}^{N-2}\left(N-1-n\right)\cdot 2^{n}
  =\sum_{k=1}^{N-1}k\cdot 2^{N-1-k}=T_{N-1},
\end{equation*}
und mit dem Lemma~\ref{lem:summenidentitaet} ($m=N-1$) ist
$T_{N-1}=2^{N}-N-1$, also
\begin{equation*}
  \left(N+1\right)+T_{N-1}
  =\left(N+1\right)+\left(2^{N}-N-1\right)=2^{N}.\qedhere
\end{equation*}
\end{proof}

\begin{bem}[Gesamtlänge]\label{bem:gesamtlaenge}
Die rechte Seite von (\ref{eqn:abgeltungsidentitaet}) zerlegt die
Bit-Länge eines Vertreters von $B_{N+1}$ nach Beiträgen der einzelnen
Läufe: Der einzige Lauf der Länge $N+1$ trägt $N+1$ Bits bei, ein Lauf
der Länge $k\leq N-1$ kommt $2^{N-1-k}$-mal vor und trägt jeweils $k$
Bits bei; die Länge $N$ trägt nichts bei, denn $_{N+1}g_N=0$
(Bemerkung~\ref{bem:keinlaufnminus1}). Das Lemma besagt also: Jeder
Vertreter von $B_{N+1}$ besitzt genau $2^{N}$ Bits, jeder Vertreter von
$B_N$ dementsprechend genau $2^{N-1}$ Bits.
\end{bem}

%-----------------------------------------------------
% 7  Algorithmen
%-----------------------------------------------------
\section{Algorithmen zur Erzeugung eines Elementes aus
\texorpdfstring{$B_N$}{B\_N}}\label{sec:algorithmen}

\noindent Motiviert durch die Ergebnisse aus dem vorherigen Abschnitt
geben wir einen iterativen Algorithmus zur Berechnung eines
Klassenvertreters der $B_N$ mit $N<\infty$ an.

\begin{figure}[htbp]
\begin{lstlisting}[caption={Iterative Erzeugung einer Folge aus B\_N fuer ein beliebiges N},label={algo:iterativ}]
Eingabe: N, die Laenge des laengsten Laufs
Lege einen leeren Bitvektor B an
Haenge einen Lauf der Laenge N an
fuer alle i von N-2 bis 1
    haenge 2^(N-i-2) Laeufe der Laenge i an B an
Gib B aus
\end{lstlisting}
\end{figure}

\noindent Die Richtigkeit des Algorithmus ergibt sich unmittelbar aus
Korollar~\ref{kor:umkehrung}: Er erzeugt genau die Laufverteilung
(\ref{eqn:grundform}), und jede Bitfolge mit dieser Laufverteilung
liegt in $B_N$.

\noindent Für den rekursiven Algorithmus nutzen wir die folgende
Konstruktion, die den Übergang von $B_{N-1}$ nach $B_N$ beschreibt.

\begin{lemma}[Konstruktion]\label{lem:konstruktion}
Sei $N\geq 3$ und $({}_{N-1}b_m)\in B_{N-1}$. Erweitere jeden Lauf von
$({}_{N-1}b_m)$ um genau ein Bit und füge $2^{N-3}$ Läufe der Länge $1$
an Stellen ein, an denen der alternierende Aufbau erhalten bleibt (etwa
je einen zwischen benachbarte Läufe und einen an ein Ende). Dann liegt
das Ergebnis in $B_N$.
\end{lemma}

\begin{proof}
Nach Satz~\ref{satz:grundverteilung}, angewandt auf $B_{N-1}$, besitzt
$({}_{N-1}b_m)$
\begin{equation*}
  \sum_{n=1}^{N-3}2^{N-3-n}+1=\left(2^{N-3}-1\right)+1=2^{N-3}
\end{equation*}
Läufe. Nach der Erweiterung besitzt die Folge die Laufängen $k+1$
anstelle von $k$ sowie $2^{N-3}$ neue Läufe der Länge $1$. Ein Lauf
der Länge $k$ enthält $k-n+1$ Teilläufe der Länge $n$; nach der
Erweiterung also $(k+1)-n+1=k-(n-1)+1$. Für die kumulierte Verteilung
$L'$ der erweiterten Folge gilt daher für alle $2\leq n\leq N$
\begin{equation*}
  l'_n=\sum_{k\geq n-1}\left(k-(n-1)+1\right)\cdot\,_{N-1}g_k
  =\,_{N-1}l_{n-1}=2^{(N-1)-(n-1)}=2^{N-n},
\end{equation*}
die Erweiterung verschiebt also die kumulierte Verteilung um eine
Stelle. Die neuen Läufe der Länge $1$ verändern nur $l'_1$: Nach
Bemerkung~\ref{bem:gesamtlaenge} hat $({}_{N-1}b_m)$ $2^{N-2}$ Bits,
die Erweiterung fügt $2^{N-3}$ Bits hinzu, die neuen Läufe weitere
$2^{N-3}$ Bits, also ist
\begin{equation*}
  l'_1=2^{N-2}+2^{N-3}+2^{N-3}=2^{N-1}=2^{N-1}.
\end{equation*}
Der längste Lauf hat die Länge $N$ (Erweiterung des Laufs der Länge
$N-1$), und es gilt $l'_n=2^{N-n}$ für alle $1\leq n\leq N$; das
Ergebnis liegt also in $B_N$.
\end{proof}

\noindent Am Beispiel $(0001)\in B_3$: Die Erweiterung aller Läufe
ergibt $(0000\,11)$ mit Läufen der Längen $4$ und $2$; die Ergänzung
von $2^{4-3}=2$ Läufen der Länge $1$, etwa durch Anhängen zu
$(0000\,11\,0\,1)$, liefert die Grundverteilung $(2;1;0;1)$ von $B_4$.

\noindent Das Lemma~\ref{lem:konstruktion} begründet die Richtigkeit
des folgenden rekursiven Algorithmus zur Berechnung eines
Klassenvertreters der $B_N$ mit $N<\infty$.

\begin{figure}[htbp]
\begin{lstlisting}[caption={Rekursive Berechnung einer Bitfolge aus einem B\_N},label={algo:rekursiv}]
Eingabe: Ein Element aus B_N mit N > 2
wiederhole fuer jeden Lauf
    ermittle die Laenge des Laufs und erweitere ihn um 1 Bit
    falls der Lauf die Laenge 1 hatte, fuege 2 Laeufe der Laenge 1 hinzu
Gib das Ergebnis aus
\end{lstlisting}
\end{figure}

Die Fälle $N=1$ und $N=2$ können in einer Implementierung des
Algorithmus gesondert behandelt werden.
% TODO(Inhalt): Referenzimplementierung (getesteter Code) veroeffentlichen
% und die Korrektheit der Algorithmen als Lemma/Laufsatz fassen.

%-----------------------------------------------------
% 8  Zusammengesetzte Verteilungen
%-----------------------------------------------------
\section{Zusammengesetzte Elemente aus verschiedenen
\texorpdfstring{$B_N$}{B\_N}}\label{sec:zusammengesetzt}

Sei nun $B\in\mathbb N$ beliebig mit seiner eindeutig bestimmten,
binären Darstellung
\begin{equation}\label{eqn:bindarstellung}
  B = \sum_{n=1}^{N} b_n\cdot 2^{n-1},\quad b_n\in\{0,1\},\ b_N=1,
\end{equation}
wobei $b_1$ das Bit mit dem kleinsten Stellenwert ist. Nach
Bemerkung~\ref{bem:gesamtlaenge} besitzt jeder Vertreter von $B_n$
genau $2^{n-1}$ Bits. Ordnen wir also dem Bit $b_n$ die Klasse $B_n$ zu,
so entspricht die Zahl $B$ genau der Zusammensetzung der Klassen
$\{B_n:\ b_n=1\}$, und die zusammengesetzte Bitfolge hat die durch $B$
vorgegebene Länge:
\begin{equation}\label{eqn:bsummelaufehaeufigkeit}
  B = \sum_{n=1}^{N} b_n\cdot 2^{n-1}
  = \sum_{\substack{1\leq n\leq N\\ b_n=1}}
  \left|\text{Vertreter von } B_n\right|.
\end{equation}
Damit können wir auf eindeutige Weise eine beliebige natürliche Zahl
$B$ mit einer durch diese Zahl eindeutig bestimmten Zusammensetzung von
zueinander fremden Häufigkeitsverteilungsklassen $B_N$ aus
$\mathcal{B_Z}$ identifizieren und erhalten die zusammengesetzte
Häufigkeitsverteilung $_BH_N$ der Läufe einer Bitfolge mit der
vorgegebenen Länge $B$. Wir vereinbaren die folgende Definition.

\begin{defi}[Zufällige Bitfolge]\label{defi:zufaellig}
Jede Bitfolge $(b_n)$, deren Häufigkeitsverteilung ihrer Läufe eine
beliebige Zusammensetzung von Häufigkeitsverteilungen zueinander
fremder Häufigkeitsverteilungsklassen $B_N \in \mathcal{B_Z}$ ist,
nennen wir zufällig.
\end{defi}
% TODO(Diskussion): Diese Definition setzt den Zufallsbegriff, der
% Gegenstand der Untersuchung ist, per Festsetzung. In Abschnitt
% \ref{sec:diskussion} gegen Martin-Lof/Kolmogoroff abgrenzen und als
% Definitionsvorschlag kennzeichnen oder begruenden.

Die Häufigkeit der Lauflänge $n$ einer beliebigen Zusammensetzung
zueinander fremder Häufigkeitsverteilungsklassen aus $\mathcal{B_Z}$
ergibt sich aus Satz~\ref{satz:grundverteilung}: Die Klasse $B_n$ trägt
(für $b_n=1$) einen einzigen Lauf der Länge $n$ bei, die Klasse
$B_{n+1}$ keinen (denn $_{n+1}g_n=0$), und jede Klasse $B_i$ mit
$i\geq n+2$ trägt $2^{\,i-n-2}$ Läufe der Länge $n$ bei:
\begin{equation}\label{eqn:hnformel}
  _BH_n=b_n+\sum_{i=n+2}^{N} b_{i}\cdot 2^{\,i-n-2},\qquad 1\leq n\leq N.
\end{equation}
Die Länge der auf diese Art gewonnenen Bitfolge ist damit
\begin{align}
  \sum_{n=1}^{N} n\cdot _BH_n
  &= \sum_{n=1}^{N} n\cdot b_n+\sum_{n=1}^{N}\;\sum_{i=n+2}^{N}
     n\cdot b_{i}\cdot 2^{\,i-n-2}\notag\\
  &= \sum_{i=1}^{N} b_i\left(i+\sum_{n=1}^{i-2} n\cdot 2^{\,i-n-2}\right)
   = \sum_{i=1}^{N} b_i\left(i+T_{i-2}\right)\notag\\
  &= \sum_{i=1}^{N} b_i\cdot 2^{\,i-1}=B,\label{eqn:laengekonsistenz}
\end{align}
mit der Summenidentität (Lemma~\ref{lem:summenidentitaet}, $m=i-2$) im
vorletzten Schritt, denn $T_{i-2}=2^{i-1}-i$. Die Zusammensetzung ist
damit konsistent mit der binären Darstellung
(\ref{eqn:bindarstellung}). Wir illustrieren diesen Sachverhalt mit
\begin{equation}
  B=101010101_2=341_{10}
\end{equation}
in der folgenden Tabelle:
\begin{table}[htbp]
\centering
\begin{tabular}{rcccccccccc}
\toprule
$n$&$b_n$&$1$&$2$&$3$&$4$&$5$&$6$&$7$&$8$&$9$\\\midrule
$1$&$1$&$ 1$&  &&  &&&& \\
$2$&$0$&$\xout    0$&$\xout  1$&   &&&&& \\
$3$&$1$&$  1$&$  0$&$  1$&&&&& \\
$4$&$0$&$\xout  2$&$\xout  1$&$\xout  0$&$\xout 1$&&&& \\
$5$&$1$&$  4$&$  2$&$  1$&$0$&$1$&&& \\
$6$&$0$&$\xout  8$&$\xout   4$&$\xout  2$&$\xout 1$&$\xout 0$&$\xout 1$&& \\
$7$&$1$&$16$&$  8$&$  4$&$2$&$1$&$0$&$1$& \\
$8$&$0$&$\xout {32}$&$\xout {16}$&$\xout   8$&$\xout 4$&$\xout 2$&$\xout 1$&$\xout 0$&$\xout 1$& \\
$9$&$1$&$64$&$32$&$16$&$8$&$4$&$2$&$1$&$0$&$1$\\\midrule
&$\Sigma$&$86$&$42$&$22$&$10$&$6$&$2$&$2$&$0$&$ 1$\\
&$n$&$1$&$2$&$3$&$4$&$5$&$6$&$7$&$8$&$9$\\
&$\Pi$&$86$&$84$&$66$&$40$&$30$&$12$&$14$&$0$&$9$\\
$\Sigma$&$341$&&&&&&&&&\\
\bottomrule
\end{tabular}
\caption{Zusammengesetzte Verteilung für $B=341$: Spaltenweise
Beiträge der Klassen $B_i$ mit $b_i=1$ zur Häufigkeit $H_n$ der
Läufe der Länge $n$; Zeile $\Sigma$: $H_n$, Zeile $\Pi$:
$n\cdot H_n$. Gestrichene Einträge gehören zu Klassen mit $b_i=0$.}
\label{tab:zusammengesetzt}
\end{table}

\noindent Wir sind nun im Stande, jede beliebige natürliche Zahl einer
endlichen, eindeutig bestimmten Teilmenge
$\{B_n\}_{n\in\mathbb N} \subset \mathcal{B_Z}$ mit
$n_1\neq n_2\ \forall n_1,n_2$ zuzuordnen und damit eine eindeutig
bestimmte, zusammengesetzte Häufigkeitsverteilung von Läufen
herzuleiten.

\noindent Wir untersuchen nun den logarithmischen Wert der Häufigkeiten
der Läufe eines solchen Klassensystems.
% TODO(Inhalt): Notation (_nK) ist nicht wohldefiniert; Abschnitt
% formal ueberarbeiten. Logarithmen mit Basis 2 kennzeichnen.
Sei $(b_n)_{n<N<\infty}$ eine Folge von binären Werten mit $N$
beliebig und fest, $b_n \in \{0,1\}$ und $b_N \neq 0$. Sei weiterhin
$_nK$ die kumulierte Häufigkeit der zu einem $n$ zugehörigen Klasse
$B_n$ aus $\mathcal{B_Z}$, also
\begin{equation}
  (_nK) =\left( 2^n, 2^{n-1}, \cdots, 2, 1\right)
\end{equation}
Sind alle $b_n=0$, $n<N$, so ist
\begin{equation}
  (_nK) _{n\leq N} =\left( 2^N, 2^{N-1}, \cdots, 2, 1\right)
\end{equation}
und die logarithmischen Werte der zusammengesetzten kumulierten
Häufigkeit der zu $(b_N)$ zugehörigen Klassen
\begin{equation}
  \log_2\left(_nK\right)_{n\leq N} =\left(N,N-1,N-2,\cdots, 2,1,0\right)
\end{equation}
Sind dagegen alle $b_n=1$, $n<N$, so sind
\begin{equation}
  (_nK) _{n\leq N} =\left(\sum_{i=1}^{N} 2^{i},\sum_{i=1}^{N-1} 2^{i}, \cdots , 1 \right)
\end{equation}
und damit
\begin{equation}
  _nK < 2^{n+1}\quad\forall n\leq N.
\end{equation}
Demnach gilt für die logarithmischen Werte der zusammengesetzten
kumulierten Häufigkeiten
\begin{equation}
  \log_2\left(_nK\right) < n+1\quad\forall n\leq N.
\end{equation}
Wechseln also die Werte $b_n$ von $0$ auf $1$, so wächst der Zweierlogarithmus
der kumulierten Häufigkeiten um maximal $1$, im Einklang mit den
experimentellen Beobachtungen aus Abschnitt~\ref{sec:experiment}.

%-----------------------------------------------------
% 9  Diskussion
%-----------------------------------------------------
\section{Diskussion}\label{sec:diskussion}

Allein die Häufigkeitsverteilung der Läufe einer Bitfolge entscheidet
über ihre Zugehörigkeit zu einem $B_N$ mit $N<\infty$. Dies hat zur
Folge, dass z.\,B.
\begin{equation}
  11110010
\end{equation}
bereits ein Element aus einem $B_N$ mit $N<\infty$ ist, mit abfallender
Länge ihrer Läufe. Als weitere Illustration können wir jede Bitfolge
als binäre Zahl auffassen und mit jeder anderen Bitfolge durch die
gewohnten Ordnungsrelationen vergleichen. Wir nehmen die Elemente aus
$B_5$. Das größte Element ist\footnote{Die Abstände in der Darstellung
dienen lediglich der Lesbarkeit und haben keinerlei Bedeutung für den
Inhalt.}
\begin{equation}
  11111\ 0\ 111\ 0\ 11\ 0\ 11\ 0
\end{equation}
Man erkennt sofort, dass die Läufe mit einer Länge größer eins in
abfallender Reihe sortiert sind und jeweils von einem Lauf der Länge
eins getrennt sind. Das kleinste Element aus $B_5$ ist dagegen
\begin{equation}
  00000\ 1\ 000\ 1\ 00\ 1\ 00\ 1
\end{equation}
Man erkennt sofort, dass es sich um das binäre Komplement des größten
Elementes aus $B_5$ handelt. Weiter ist
\begin{equation}\label{eqn:b5kleinstemitseins}
  1\ 00000\ 1\ 000\ 1\ 00\ 1\ 00
\end{equation}
das kleinste Element aus $B_5$, das mit einer Eins anfängt, und sein
binäres Komplement
\begin{equation}\label{eqn:b5groesstermitnull}
  0\ 11111\ 0\ 111\ 0\ 11\ 0\ 11
\end{equation}
das größte Element, das mit einer Null anfängt. Weiterhin gibt es in
$B_5$ kein Element zwischen (\ref{eqn:b5kleinstemitseins}) und
(\ref{eqn:b5groesstermitnull}).

\noindent So lässt sich intuitiv eine gewisse Struktur in den $B_N$
erahnen. Es kommt nun nur noch darauf an, dass diese
Wiedererkennbarkeit methodologisch nicht relevant ist, um die $B_N$ als
Menge der zufälligen Bitfolgen zu bezeichnen. Vielmehr ist hier
angemessen zu behaupten, dass es den Zufall, wie er in der
Wahrscheinlichkeitstheorie begriffen wird, nicht gibt, sondern die
Willkür der Wahl eines Elementes aus einem beliebigen $B_N$, welche nur
mit exponentiellem Aufwand zurückverfolgt werden kann. Dieser
asymmetrische Aufwand einer Entscheidung in (super-)linearer Zeit und
einer Rückverfolgung in exponentieller Zeit ist Grundlage der
Kryptographie und gilt, solange
\begin{equation}
  P \neq NP.
\end{equation}
% TODO(Inhalt, kritisch): Die P!=NP-Aussage ist unbelegt und in dieser
% Form nicht haltbar: Es ist weder eine Praezisierung des "Rueckverfolgungs-
% aufwands" gegeben, noch ist die behauptete Verbindung zu P vs. NP
% hergeleitet. Entweder als praezise, beweisbare Aussage ueber die
% Klassenvertreter formulieren oder streichen. Ebenso die These "es gibt
% den Zufall nicht" als philosophische Behauptung kennzeichnen und an
% Popper/Martin-Lof zurueckbinden.

%-----------------------------------------------------
% 10  Zusammenfassung und Ausblick
%-----------------------------------------------------
\section{Zusammenfassung und Ausblick}\label{sec:schluss}

Wir haben mit dem Einschritttest ein Auswertungsverfahren für endliche
Bitfolgen eingeführt, das die Häufigkeiten von Läufen und Teilläufen
kumuliert. Die experimentelle Anwendung auf die Ausgabe von
\texttt{/dev/random} zeigt einen exponentiell fallenden Verlauf, wie er
aus der geometrischen Verteilung der Lauflängen unabhängiger Bits zu
erwarten ist. Die Klassen $B_N$ der Bitfolgen mit exakt exponentiell
fallender kumulierter Verteilung besitzen eine streng rekursive
Grundverteilung (Satz~\ref{satz:grundverteilung}); daraus ergeben sich
ein iterativer und ein rekursiver Generator, und über die binäre
Darstellung natürlicher Zahlen eine eindeutige Zusammensetzung der
Klassen (Abschnitt~\ref{sec:zusammengesetzt}).

Offen und Gegenstand der Weiterarbeit sind:
\begin{itemize}
  \item die Neu-Erhebung und Veröffentlichung der Versuchsdaten
        (exakte Bitanzahl, Erfassungsskript, Korrektur der
        Ausreißer in Tabelle~\ref{tab:haeufigkeiten}) sowie eine
        statistische Absicherung gegen die theoretische Verteilung
        (Anpassungstest);
  \item eine Referenzimplementierung der Algorithmen
        \ref{algo:iterativ} und \ref{algo:rekursiv} samt Test
        der Klassenmitgliedschaft der erzeugten Folgen;
  \item die Abgrenzung des in Definition~\ref{defi:zufaellig}
        vorgeschlagenen Zufallsbegriffs gegenüber dem
        wahrscheinlichkeitstheoretischen und dem algorithmischen
        Zufallsbegriff \cite{kolmogoroff1933,martinloef1966};
  \item eine Präzisierung der in Abschnitt~\ref{sec:diskussion}
        angedeuteten Komplexitätsaussagen.
\end{itemize}

%-----------------------------------------------------
% Literatur
%-----------------------------------------------------
\printbibliography[title={Literaturverzeichnis}]

\end{document}
